% GENERATED FILE. Edit canonical lessons, then run npm run generate:books.
% canonical-source-hash: fnv1a64:99d91a78cf7722bd
% canonical-lessons: KA-C279-cuccu, KA-C279-ugulu, KA-C279-akalisu, KA-C279-bageharisu, KA-C279-age

\chapter{To prick, To spit, To yawn, To solve, To dig}
\label{ch:ka-to-prick-to-spit-to-yawn-to-solve-to-dig}
\hlchaptermodality{279}

\begin{chapteropening}
\textbf{By the end of this chapter:} \emph{I can say to prick, to spit, to yawn, to solve and to dig.}

Complete the last lesson of chapter 279: I can say to prick, to spit, to yawn, to solve and to dig.
\end{chapteropening}

\section[\texorpdfstring{\kn{ಚುಚ್ಚು} (cuccu)}{ಚುಚ್ಚು (cuccu)}]{\kn{ಚುಚ್ಚು} (cuccu) — to prick, to pierce}

\label{lesson:KA-C279-cuccu}



\begin{quote}
Before the new one: say the Kannada for to prevent, then the Kannada for to prepare.
\end{quote}

\subsection*{You'll want to know: \kn{ಚುಚ್ಚು}}
\textbf{\kn{ಚುಚ್ಚು}} — \emph{cuccu} — \textquotedblleft{}to prick, to pierce\textquotedblright{}.

\begin{grammarlens}[title={What you've built}]
One more word for this chapter.
\end{grammarlens}

\subsection*{Your turn}
\begin{itemize}
\raggedright
  \item \emph{Say it:} \emph{cuccu}
  \item \emph{Say it:} \emph{cuccu}, once more
  \item \emph{Say it:} say \emph{cuccu}
  \item \emph{Recall:} say the Kannada for to prevent, then the Kannada for to prepare, then say \emph{cuccu} again
\end{itemize}

\begin{tcolorbox}[breakable,colback=teal!4,colframe=teal!35!black,title={Before you move on}]
What does \kn{ಚುಚ್ಚು} mean? (\textquotedblleft{}To prick, to pierce\textquotedblright{}.) Say it once more.
\end{tcolorbox}
\section[\texorpdfstring{\kn{ಉಗುಳು} (uguḷu)}{ಉಗುಳು (uguḷu)}]{\kn{ಉಗುಳು} (uguḷu) — to spit}

\label{lesson:KA-C279-ugulu}



\begin{quote}
Before the new one: say the Kannada for to prepare, then the Kannada for to prick.
\end{quote}

\subsection*{You'll want to know: \kn{ಉಗುಳು}}
\textbf{\kn{ಉಗುಳು}} — \emph{uguḷu} — \textquotedblleft{}to spit\textquotedblright{}.

\begin{grammarlens}[title={What you've built}]
One more word for this chapter.
\end{grammarlens}

\subsection*{Your turn}
\begin{itemize}
\raggedright
  \item \emph{Say it:} \emph{uguḷu}
  \item \emph{Say it:} \emph{uguḷu}, once more
  \item \emph{Say it:} say \emph{uguḷu}
  \item \emph{Recall:} say the Kannada for to prepare, then the Kannada for to prick, then say \emph{uguḷu} again
\end{itemize}

\begin{tcolorbox}[breakable,colback=teal!4,colframe=teal!35!black,title={Before you move on}]
What does \kn{ಉಗುಳು} mean? (\textquotedblleft{}To spit\textquotedblright{}.) Say it once more.
\end{tcolorbox}
\section[\texorpdfstring{\kn{ಆಕಳಿಸು} (ākaḷisu)}{ಆಕಳಿಸು (ākaḷisu)}]{\kn{ಆಕಳಿಸು} (ākaḷisu) — to yawn}

\label{lesson:KA-C279-akalisu}



\begin{quote}
Before the new one: say the Kannada for to prick, then the Kannada for to spit.
\end{quote}

\subsection*{You'll want to know: \kn{ಆಕಳಿಸು}}
\textbf{\kn{ಆಕಳಿಸು}} — \emph{ākaḷisu} — \textquotedblleft{}to yawn\textquotedblright{}.

\begin{grammarlens}[title={What you've built}]
One more word for this chapter.
\end{grammarlens}

\subsection*{Your turn}
\begin{itemize}
\raggedright
  \item \emph{Say it:} \emph{ākaḷisu}
  \item \emph{Say it:} \emph{ākaḷisu}, once more
  \item \emph{Say it:} say \emph{ākaḷisu}
  \item \emph{Recall:} say the Kannada for to prick, then the Kannada for to spit, then say \emph{ākaḷisu} again
\end{itemize}

\begin{tcolorbox}[breakable,colback=teal!4,colframe=teal!35!black,title={Before you move on}]
What does \kn{ಆಕಳಿಸು} mean? (\textquotedblleft{}To yawn\textquotedblright{}.) Say it once more.
\end{tcolorbox}
\section[\texorpdfstring{\kn{ಬಗೆಹರಿಸು} (bageharisu)}{ಬಗೆಹರಿಸು (bageharisu)}]{\kn{ಬಗೆಹರಿಸು} (bageharisu) — to solve, to sort out}

\label{lesson:KA-C279-bageharisu}



\begin{quote}
Before the new one: say the Kannada for to spit, then the Kannada for to yawn.
\end{quote}

\subsection*{You'll want to know: \kn{ಬಗೆಹರಿಸು}}
\textbf{\kn{ಬಗೆಹರಿಸು}} — \emph{bageharisu} — \textquotedblleft{}to solve, to sort out\textquotedblright{}.

\begin{grammarlens}[title={What you've built}]
One more word for this chapter.
\end{grammarlens}

\subsection*{Your turn}
\begin{itemize}
\raggedright
  \item \emph{Say it:} \emph{bageharisu}
  \item \emph{Say it:} \emph{bageharisu}, once more
  \item \emph{Say it:} say \emph{bageharisu}
  \item \emph{Recall:} say the Kannada for to spit, then the Kannada for to yawn, then say \emph{bageharisu} again
\end{itemize}

\begin{tcolorbox}[breakable,colback=teal!4,colframe=teal!35!black,title={Before you move on}]
What does \kn{ಬಗೆಹರಿಸು} mean? (\textquotedblleft{}To solve, to sort out\textquotedblright{}.) Say it once more.
\end{tcolorbox}
\section[\texorpdfstring{\kn{ಅಗೆ} (age)}{ಅಗೆ (age)}]{\kn{ಅಗೆ} (age) — to dig}

\label{lesson:KA-C279-age}



\begin{quote}
Before the new one: say the Kannada for to yawn, then the Kannada for to solve.
\end{quote}

\subsection*{You'll want to know: \kn{ಅಗೆ}}
\textbf{\kn{ಅಗೆ}} — \emph{age} — \textquotedblleft{}to dig\textquotedblright{}.

\begin{grammarlens}[title={What you've built}]
One more word for this chapter.
\end{grammarlens}

\subsection*{Your turn}
\begin{itemize}
\raggedright
  \item \emph{Say it:} \emph{age}
  \item \emph{Say it:} \emph{age}, once more
  \item \emph{Say it:} say \emph{age}
  \item \emph{Recall:} say the Kannada for to yawn, then the Kannada for to solve, then say \emph{age} again
\end{itemize}

\begin{tcolorbox}[breakable,colback=teal!4,colframe=teal!35!black,title={Before you move on}]
What does \kn{ಅಗೆ} mean? (\textquotedblleft{}To dig\textquotedblright{}.) Say it once more.
\end{tcolorbox}
